% !TeX root = ../papers/nc-p3.tex
\section{余因子曲線の整性によらない二つの表示の排除}
\label{sec:unrestricted-binary-formats}
\begin{lemma}\label{lem:finite-nine-on-plane}
四辺有限被約な標準有限Frobenius窓の各位置では、同一平面に載る印は九以下である。
その入力を共有する別の有限六双一次完全交叉辺にも純テンソル関係はない。
\end{lemma}
\begin{proof}
三次評価同型から平面に載る印は高々十。十個の集合 $A$ が平面にあると仮定し、補集合を $B$ とする。
隣接位置を $U$ と書く。$X^2,U^2$ は非零対角重み付きの完全直交補空間なので、
$B$ での $U^2$ の制限階数は $|B|-10+\operatorname{rank}(X^2|_A)\le6$。
もし $U$ の $B$ の十点が非共面なら、その座標環のHilbert関数は一次で四、二次で六以下、
三次で十となる（全二十点で三次評価が同型）。一般一次非零因子で割ると二次の次元は二以下であり、
Macaulayの増大上界により三次も二以下である。従って元の三次次元は八以下で矛盾する。
二平面の式の積は有限グラフ上の純関係となり、標準有限Frobenius窓の有限辺での純関係の不在に反する。
後半では別の辺の相手射影でも三次評価が同型で平面切断は高々十であることを使う。
純関係なら最初の平面の外の少なくとも十一印が他方の平面に載ってしまう。
\end{proof}

\begin{theorem}\label{thm:unrestricted-bezout-exclusion}
被約四辺有限標準正則窓の非恒等な反対根は、入出力の基底を除いて
$\operatorname{Bez}(T(t_0x),T(t_1x))$, $T\in\operatorname{GL}(S_4)$ という形を持たない。
余因子曲線の成分数・一般重複度、階数一点の有限性、行列式の整性を仮定しない。
特に整な行列式八次では六次余因子の包は九次元または十次元である。
\end{theorem}
\begin{proof}
基底変更行列 $C$ の空間へ射影した標準辺の退化スキームは、$T$ の可逆性だけで有限で長さ二十である
（系\ref{cor:trisecant-middle-edge}）。この点での六行四列の階数を $r$ とする。
$r\le1$ なら繊維の平面上の全 $x$ で三次共通因子があり、根行列の階数は一以下。
階数二以下の曲面の排除（定理\ref{thm:rank-two-surface-excess-curves}）に反する。
$r=2$ なら繊維は直線であり、実際の $X$ から来る印は高々三。しかしこの $C$ の局所退化スキームは
長さ四以上である。二行を消去した最大小行列式は全て極大イデアルの平方に入るためである。
$r=3$ なら局所長は一以上で $X$ の印は高々一。
全ての $X$ の印は一意な $C$ を持つので、二十印と総長二十を比較すると階数二点はなく、
全ての局所長は一である。従って元の三次共通因子軌跡も有限で標準辺は両側有限被約である。

もし二つの四次評価曲線に交点があれば、ある $z,w$ で $T^*\mathrm{ev}_z$ が
$\mathrm{ev}_w$ に比例する。標準辺の関係を $z$ で評価すれば
$x(w)$ と $C$ の一つの非零一次式の積となり、前補題に反する。
従って正規化候補の母線束は全底上で定義され、入力射は有限である。
その像は結果式の台そのもので整であり、一般に二つ以上の共通根（長さで数える）があれば
像曲面上で根行列の階数二以下となり、再び曲面排除に反する。
従って射は双有理で、母線束の超平面自乗八から結果式は重複しない整八次である。
この状況で定理\ref{thm:exclude-integral-cofactor-curves}の母線束の証明は
余因子曲線の整性をもはや用いない。六十の繊維長、二十印の外での繊維長二以下、
二つの四次評価曲線の不一致をいま全て確認したので、同じ二つの商線束から正則性に矛盾する。
整八次での最後の主張は定理\ref{thm:seven-cofactor-trisecant-grassmannian}と八余因子を排除する定理\ref{thm:eight-cofactor-excess}による。
\end{proof}

\begin{theorem}\label{thm:finite-quintic-web-exclusion}
被約四辺有限標準正則窓のある位置の二十印が、五行対称一次行列 $A(x)$ の
階数三以下スキームそのもので、$\det A$ が非零であるとする。
この位置と反対位置の標識付き配置は射影同値である。五次式の既約性も被約性も仮定しない。
特に二次曲面正規化の表示 $\nu_*\mathcal O_Q(1)$ で長さ三繊維の軌跡が有限なら、
余因子曲線が可約・肥厚でも非恒等根として正則窓に入らない。
\end{theorem}
\begin{proof}
各印の階数は三である。二以下なら四次小行列式が極大イデアルの平方に入り、被約点になれない。
また四次小行列式の対称分解から $I_X(4)$ の全基底を得る。
定理\ref{thm:quadric-conductor-quintic}の四次対合の証明で、
五次 $F$ が既約であることを用いた箇所をスキームとして置き換える。
入射 $A(x)p=0$ は階数四開集合で五次零スキームに同型、階数三の有限集合では次元一の繊維を持つ。
従って $[A_i p]$ の各階数下界の余次元評価、$b(p)$ の余次元二の完全なHilbert--Burch分解、
四次双有理対合の議論はそのまま成立する。特に $Z=V(b_i)$ は埋込み成分を持たない。
各曲面成分の一般点では $[A_i p]$ の階数が三で $x$ は一意、かつ $A(x)$ の階数が四である
（固定した二十印の逆像は直線の有限和）。ここで入射は $Z$ と $F$ の両方に局所的に同型である。
$\operatorname{adj}A$ は $F$ 上の単元四行小行列式で消去すると、その局所環で
核ベクトル $p$ の平方の単元倍となる。これは非被約な $F$ でも成立する。
$\Phi(\operatorname{adj}A(x))$ のPlücker式は入力で恒等零だから、
$\Phi(p^2)$ の四行Pfaffianは $Z$ の各一般点でスキームとして消える。
純性よりそれらは $I_Z$ に入り、四次式なので五つの $b_i$ の線形結合である。
以後の四次対合、階数四の線形空間、二項四次束の証明は $F$ の整性に依存しない。
よって全ての隣接辺が同一の $A$ のHankel構造を与える。
定理\ref{thm:exclude-degree-twenty-hankel-pairs}の比較では $A(x)$ の階数下界三だけで足り、
非恒等反対根が排除される。二次曲面に関する主張では、正規化環の六行対称表示と単元生成元の議論により
$A$ の階数三以下が長さ三繊維の有限集合で、対称分解と既知の二十印からそのスキームは被約である。
\end{proof}
\path{computations/verify_unrestricted_binary_formats.py} は平面補集合・局所長の境界と、
重複因子を持つ局所対称行列での余因子平方を検算する。
